% Part: incompleteness
% Chapter: representability-in-q
% Section: basic-representable

\documentclass[../../../include/open-logic-section]{subfiles}

\begin{document}

\olfileid{inc}{req}{bre}
\olsection{الدوال الأساسية قابلة للتمثيل في~$\Th{Q}$}

أول ما يلزمنا قابلية الدوال الأساسية جميعًا للتمثيل في~$\Th{Q}$.
ومعنى المطلوب أن نعيّن !!a{formula}
$!A_\OLId{latin-ordinary}{f}(\OLId{latin-ordinary}{x}_\OLMathNumeral{0},\dots,\OLId{latin-ordinary}{x}_{\OLId{latin-ordinary}{k}-\OLMathNumeral{1}},\OLId{latin-ordinary}{y})$ تمثل، لكل عدد وسائط $\OLId{latin-ordinary}{k}$،
الدالة الأساسية $\OLId{latin-ordinary}{f}(\OLId{latin-ordinary}{x}_\OLMathNumeral{0},\dots,\OLId{latin-ordinary}{x}_{\OLId{latin-ordinary}{k}-\OLMathNumeral{1}})$.

وسنشمل بذلك الصفر والخلف والجمع والضرب والدالة المميزة للمساواة
والإسقاطات. واختيار الصيغة في هذه المواضع قريب: للصفر الصيغة $\OLId{latin-ordinary}{y}=\Obj 0$،
وللخلف !!{formula} $\OLId{latin-ordinary}{x}_\OLMathNumeral{0}'=\OLId{latin-ordinary}{y}$، وللجمع !!{formula} $(\OLId{latin-ordinary}{x}_\OLMathNumeral{0}+\OLId{latin-ordinary}{x}_\OLMathNumeral{1})=\OLId{latin-ordinary}{y}$.
وإنما العمل في إثبات أن $\Th{Q}$ تبرهن ما يلزم من !!{sentence}s.
فالقول، مثلًا، إن هذه الـ!!{formula} تمثل الجمع معناه أن $\Th{Q}$
تثبت، لكل عددين طبيعيين $\OLId{latin-ordinary}{m}$ و$\OLId{latin-ordinary}{n}$، ما يأتي:
\begin{align*}
& \eq[\num \OLId{latin-ordinary}{n} + \num \OLId{latin-ordinary}{m}][\num {\OLId{latin-ordinary}{n}+\OLId{latin-ordinary}{m}}] \text{ و}\\
& \lforall[\OLId{latin-ordinary}{y}][(\eq[(\num \OLId{latin-ordinary}{n} + \num \OLId{latin-ordinary}{m})][\OLId{latin-ordinary}{y}] \lif \eq[\OLId{latin-ordinary}{y}][\num{\OLId{latin-ordinary}{n}+\OLId{latin-ordinary}{m}}])].
\end{align*}

\begin{prop}
\ollabel{prop:rep-zero}
الصيغة $!A_{\Zero}(\OLId{latin-ordinary}{x},\OLId{latin-ordinary}{y})\ident\eq[\OLId{latin-ordinary}{y}][\Obj 0]$ تمثيل لدالة الصفر
$\Zero(\OLId{latin-ordinary}{x})=\OLMathNumeral{0}$ في~$\Th{Q}$.
\end{prop}

\begin{prop}
\ollabel{prop:rep-succ}
الصيغة $!A_{\Succ}(\OLId{latin-ordinary}{x},\OLId{latin-ordinary}{y})\ident\eq[\OLId{latin-ordinary}{y}][\OLId{latin-ordinary}{x}']$ تمثيل لدالة الخلف
$\Succ(\OLId{latin-ordinary}{x})=\OLId{latin-ordinary}{x}+\OLMathNumeral{1}$ في~$\Th{Q}$.
\end{prop}

\begin{prop}
\ollabel{prop:rep-proj}
الصيغة
\[!A_{\Proj{\OLId{latin-ordinary}{n}}{\OLId{latin-ordinary}{i}}}(\OLId{latin-ordinary}{x}_\OLMathNumeral{0},\dots,\OLId{latin-ordinary}{x}_{\OLId{latin-ordinary}{n}-\OLMathNumeral{1}},\OLId{latin-ordinary}{y})\ident\eq[\OLId{latin-ordinary}{y}][\OLId{latin-ordinary}{x}_\OLId{latin-ordinary}{i}].\]
تمثيل لدالة الإسقاط $\Proj{\OLId{latin-ordinary}{n}}{\OLId{latin-ordinary}{i}}(\OLId{latin-ordinary}{x}_\OLMathNumeral{0},\dots,\OLId{latin-ordinary}{x}_{\OLId{latin-ordinary}{n}-\OLMathNumeral{1}})=\OLId{latin-ordinary}{x}_\OLId{latin-ordinary}{i}$ في~$\Th{Q}$.
\end{prop}

\begin{prob}
بيّن أن الصيغ $\eq[\OLId{latin-ordinary}{y}][\Obj 0]$ و$\eq[\OLId{latin-ordinary}{y}][\OLId{latin-ordinary}{x}']$ و$\eq[\OLId{latin-ordinary}{y}][\OLId{latin-ordinary}{x}_\OLId{latin-ordinary}{i}]$
تمثيلات للدوال $\Zero$ و$\Succ$ و$\Proj{\OLId{latin-ordinary}{n}}{\OLId{latin-ordinary}{i}}$، على هذا الترتيب.
\end{prob}

\begin{prop}
\ollabel{prop:rep-id}
لدالة المساواة $=$ المميزة، وهي
\[
\Char{=}(\OLId{latin-ordinary}{x}_\OLMathNumeral{0}, \OLId{latin-ordinary}{x}_\OLMathNumeral{1}) =
\begin{cases}
  \OLMathNumeral{1} & \text{إذا كان } \OLId{latin-ordinary}{x}_\OLMathNumeral{0}=\OLId{latin-ordinary}{x}_\OLMathNumeral{1}\\
  \OLMathNumeral{0} & \text{في غير ذلك}
\end{cases}
\]
تمثيل في~$\Th{Q}$ تعطيه الصيغة
\[
  !A_{\Char{=}}(\OLId{latin-ordinary}{x}_\OLMathNumeral{0}, \OLId{latin-ordinary}{x}_\OLMathNumeral{1}, \OLId{latin-ordinary}{y}) \ident (\eq[\OLId{latin-ordinary}{x}_\OLMathNumeral{0}][\OLId{latin-ordinary}{x}_\OLMathNumeral{1}] \land \eq[\OLId{latin-ordinary}{y}][\num{\OLMathNumeral{1}}]) \lor (\eq/[\OLId{latin-ordinary}{x}_\OLMathNumeral{0}][\OLId{latin-ordinary}{x}_\OLMathNumeral{1}] \land
\eq[\OLId{latin-ordinary}{y}][\num{\OLMathNumeral{0}}]).
\]
\end{prop}

ولإثبات ذلك نقدّم اللمة الآتية.

\begin{lem}
\ollabel{lem:q-proves-neq} لتكن $\OLId{latin-ordinary}{n}$ و$\OLId{latin-ordinary}{m}$ عددين طبيعيين.
فإن كان $\OLId{latin-ordinary}{n}\neq \OLId{latin-ordinary}{m}$ لزم $\Th{Q}\Proves\eq/[\num \OLId{latin-ordinary}{n}][\num \OLId{latin-ordinary}{m}]$.
\end{lem}

\begin{proof}
نجري الاستقراء على $\OLId{latin-ordinary}{n}$، والدعوى لكل $\OLId{latin-ordinary}{m}$ هي: إذا كان $\OLId{latin-ordinary}{n}\neq \OLId{latin-ordinary}{m}$،
فإن $\OLId{latin-looped}{Q}\Proves\eq/[\num \OLId{latin-ordinary}{n}][\num \OLId{latin-ordinary}{m}]$.

أما الأساس فهو $\OLId{latin-ordinary}{n}=\OLMathNumeral{0}$. فإن خالف $\OLId{latin-ordinary}{m}$ العدد $\OLMathNumeral{0}$، كتبناه
$\OLId{latin-ordinary}{m}=\OLId{latin-ordinary}{k}+\OLMathNumeral{1}$ لعدد طبيعي $\OLId{latin-ordinary}{k}$. ولدينا البديهية
$\lforall[\OLId{latin-ordinary}{x}][\eq/[\OLMathNumeral{0}][\OLId{latin-ordinary}{x}']]$؛ فنضع $\num \OLId{latin-ordinary}{k}$ موضع $\OLId{latin-ordinary}{x}$ ببديهية الكم،
ونستنتج $\eq/[\OLMathNumeral{0}][\num \OLId{latin-ordinary}{k}']$. ثم إن $\num \OLId{latin-ordinary}{k}'$ هو بعينه $\num \OLId{latin-ordinary}{m}$.

وأما خطوة الاستقراء، فنفرض الدعوى لـ$\OLId{latin-ordinary}{n}$ ثم نثبتها لـ$\OLId{latin-ordinary}{n}+\OLMathNumeral{1}$.
خذ $\OLId{latin-ordinary}{m}$ عددًا طبيعيًا كيفما كان. إما أن يكون $\OLId{latin-ordinary}{m}=\OLMathNumeral{0}$، وإما أن يكون
$\OLId{latin-ordinary}{m}=\OLId{latin-ordinary}{k}+\OLMathNumeral{1}$ لعدد $\OLId{latin-ordinary}{k}$. فالحالة الأولى على مثال حالة الأساس.
وفي الثانية، إذا فرضنا $\OLId{latin-ordinary}{n}+\OLMathNumeral{1}\neq \OLId{latin-ordinary}{k}+\OLMathNumeral{1}$، لزم $\OLId{latin-ordinary}{n}\neq \OLId{latin-ordinary}{k}$،
وبفرض الاستقراء لـ$\OLId{latin-ordinary}{n}$ نحصل على
$\Th{Q}\Proves\eq/[\num \OLId{latin-ordinary}{n}][\num \OLId{latin-ordinary}{k}]$. ومن البديهية
$\lforall[\OLId{latin-ordinary}{x}][\lforall[\OLId{latin-ordinary}{y}][\eq[\OLId{latin-ordinary}{x}'][\OLId{latin-ordinary}{y}']\lif\eq[\OLId{latin-ordinary}{x}][\OLId{latin-ordinary}{y}]]]$
نستخرج، ببديهية الكم،
$\eq[\num \OLId{latin-ordinary}{n}'][\num \OLId{latin-ordinary}{k}']\lif\eq[\num \OLId{latin-ordinary}{n}][\num \OLId{latin-ordinary}{k}]$.
ثم نستنتج بالمنطق القضي، داخل $\Th{Q}$،
$\eq/[\num \OLId{latin-ordinary}{n}][\num \OLId{latin-ordinary}{k}]\lif\eq/[\num \OLId{latin-ordinary}{n}'][\num \OLId{latin-ordinary}{k}']$.
فتعطينا قاعدة إثبات المقدّم $\eq/[\num \OLId{latin-ordinary}{n}'][\num \OLId{latin-ordinary}{k}']$.
وهذا المطلوب، لأن $\num \OLId{latin-ordinary}{k}'$ هو $\num \OLId{latin-ordinary}{m}$.
\end{proof}

\begin{explain}
ومضمون اللمة متواضع: تستطيع $\Th{Q}$ أن تثبت أن أسماء الأعداد المختلفة
تدل على أشياء مختلفة. ومن أمثلته أن تثبت $\Th{Q}$ العبارة
$\OLMathNumeral{0}''\neq\OLMathNumeral{0}'''$. غير أن تعميم هذا القول يقتضي بعض العناية.
وينبغي تمييز موضع الاستقراء هنا؛ فهو، وإن استعملناه في البرهان،
استقراء \emph{خارج} $\Th{Q}$، لا بديهية استقراء فيها.
\end{explain}

\begin{proof}[برهان \olref{prop:rep-id}]
نقسم أولًا بحسب تساوي العددين واختلافهما. فإن كان $\OLId{latin-ordinary}{n}=\OLId{latin-ordinary}{m}$ كان
$\num{\OLId{latin-ordinary}{n}}$ و$\num{\OLId{latin-ordinary}{m}}$ اسمًا واحدًا، وكانت $\Char{=}(\OLId{latin-ordinary}{n},\OLId{latin-ordinary}{m})=\OLMathNumeral{1}$.
ولما كان
$\Th{Q}\Proves(\eq[\num{\OLId{latin-ordinary}{n}}][\num{\OLId{latin-ordinary}{m}}]\land\eq[\num{\OLMathNumeral{1}}][\num{\OLMathNumeral{1}}])$،
ثبتت $!A_=(\num{\OLId{latin-ordinary}{n}},\num{\OLId{latin-ordinary}{m}},\num{\OLMathNumeral{1}})$.
وأما إذا كان $\OLId{latin-ordinary}{n}\neq \OLId{latin-ordinary}{m}$، فالقيمة $\Char=(\OLId{latin-ordinary}{n},\OLId{latin-ordinary}{m})=\OLMathNumeral{0}$.
ومن \olref{lem:q-proves-neq} نحصل على
$\Th{Q}\Proves\eq/[\num{\OLId{latin-ordinary}{n}}][\num{\OLId{latin-ordinary}{m}}]$؛ فتثبت أيضًا
$(\eq/[\num{\OLId{latin-ordinary}{n}}][\num{\OLId{latin-ordinary}{m}}]\land\Obj 0=\Obj 0)$،
ومنها $\Th{Q}\Proves !A_=(\num{\OLId{latin-ordinary}{n}},\num{\OLId{latin-ordinary}{m}},\num{\OLMathNumeral{0}})$.

ويبقى شرط انفراد القيمة، وله الحالتان نفسيهما.
إذا كان $\OLId{latin-ordinary}{n}=\OLId{latin-ordinary}{m}$، فالمطلوب
$\Th{Q}\Proves\lforall[\OLId{latin-ordinary}{y}][(!A_=(\num{\OLId{latin-ordinary}{n}},\num{\OLId{latin-ordinary}{m}},\OLId{latin-ordinary}{y})
\lif\eq[\OLId{latin-ordinary}{y}][\num{\OLMathNumeral{1}}])]$.
نبيّن الحجة بعبارة غير صورية: لنفرض
$!A_=(\num{\OLId{latin-ordinary}{n}},\num{\OLId{latin-ordinary}{m}},\OLId{latin-ordinary}{y})$، أي
\[
(\eq[\num{\OLId{latin-ordinary}{n}}][\num{\OLId{latin-ordinary}{n}}] \land \eq[\OLId{latin-ordinary}{y}][\num{\OLMathNumeral{1}}]) \lor
(\eq/[\num{\OLId{latin-ordinary}{n}}][\num{\OLId{latin-ordinary}{n}}] \land \eq[\OLId{latin-ordinary}{y}][\num{\OLMathNumeral{0}}])
\]
فالشق الأيسر يعطي $\eq[\OLId{latin-ordinary}{y}][\num{\OLMathNumeral{1}}]$ بالمنطق.
وأما الشق الأيمن فيناقض $\eq[\num{\OLId{latin-ordinary}{n}}][\num{\OLId{latin-ordinary}{n}}]$، وهي ثابتة بالمنطق.

وأما إذا كان $\OLId{latin-ordinary}{n}\neq \OLId{latin-ordinary}{m}$، فإن
$!A_=(\num{\OLId{latin-ordinary}{n}},\num{\OLId{latin-ordinary}{m}},\OLId{latin-ordinary}{y})$ هي
\[
(\eq[\num{\OLId{latin-ordinary}{n}}][\num{\OLId{latin-ordinary}{m}}] \land \eq[\OLId{latin-ordinary}{y}][\num{\OLMathNumeral{1}}]) \lor
(\eq/[\num{\OLId{latin-ordinary}{n}}][\num{\OLId{latin-ordinary}{m}}] \land \eq[\OLId{latin-ordinary}{y}][\num{\OLMathNumeral{0}}])
\]
فالشق الأيسر الآن يناقض $\eq/[\num{\OLId{latin-ordinary}{n}}][\num{\OLId{latin-ordinary}{m}}]$، وقد ثبتت في
$\Th{Q}$ باللمة \olref{lem:q-proves-neq}.
ويبقى الشق الأيمن، وهو يعطي $\eq[\OLId{latin-ordinary}{y}][\num{\OLMathNumeral{0}}]$.
\end{proof}

\begin{prop}
\ollabel{prop:rep-add}
لدالة الجمع $\Add(\OLId{latin-ordinary}{x}_\OLMathNumeral{0},\OLId{latin-ordinary}{x}_\OLMathNumeral{1})=\OLId{latin-ordinary}{x}_\OLMathNumeral{0}+\OLId{latin-ordinary}{x}_\OLMathNumeral{1}$ تمثيل في~$\Th{Q}$ بالصيغة:
\[
  !A_{\Add}(\OLId{latin-ordinary}{x}_\OLMathNumeral{0}, \OLId{latin-ordinary}{x}_\OLMathNumeral{1}, \OLId{latin-ordinary}{y}) \ident \eq[\OLId{latin-ordinary}{y}][(\OLId{latin-ordinary}{x}_\OLMathNumeral{0} + \OLId{latin-ordinary}{x}_\OLMathNumeral{1})].
\]
\end{prop}

\begin{lem}
\ollabel{lem:q-proves-add}
$\Th{Q} \Proves \eq[(\num{\OLId{latin-ordinary}{n}} + \num{\OLId{latin-ordinary}{m}})][\num{\OLId{latin-ordinary}{n}+\OLId{latin-ordinary}{m}}]$
\end{lem}

\begin{proof}
البرهان بالاستقراء على~$\OLId{latin-ordinary}{m}$. ففي الأساس $\OLId{latin-ordinary}{m}=\OLMathNumeral{0}$،
والمطلوب $\Th{Q}\Proves\eq[(\num{\OLId{latin-ordinary}{n}}+\Obj 0)][\num{\OLId{latin-ordinary}{n}}]$،
وهو نتيجة البديهية $!Q_\OLMathNumeral{4}$. ثم نفرض الدعوى لـ$\OLId{latin-ordinary}{m}$، وننتقل إلى $\OLId{latin-ordinary}{m}+\OLMathNumeral{1}$،
فيكون المطلوب
$\Th{Q}\Proves\eq[(\num{\OLId{latin-ordinary}{n}}+\num{\OLId{latin-ordinary}{m}+\OLMathNumeral{1}})][\num{\OLId{latin-ordinary}{n}+\OLId{latin-ordinary}{m}+\OLMathNumeral{1}}]$.
واسم العدد $\num{\OLId{latin-ordinary}{m}+\OLMathNumeral{1}}$ هو $\num{\OLId{latin-ordinary}{m}}'$، كما أن $\num{\OLId{latin-ordinary}{n}+\OLId{latin-ordinary}{m}+\OLMathNumeral{1}}$
هو $\num{\OLId{latin-ordinary}{n}+\OLId{latin-ordinary}{m}}'$. فبالبديهية $!Q_\OLMathNumeral{5}$ نحصل على
$\Th{Q}\Proves\eq[(\num{\OLId{latin-ordinary}{n}}+\num{\OLId{latin-ordinary}{m}}')][(\num{\OLId{latin-ordinary}{n}}+\num{\OLId{latin-ordinary}{m}})']$.
وبفرض الاستقراء لدينا
$\Th{Q}\Proves\eq[(\num{\OLId{latin-ordinary}{n}}+\num{\OLId{latin-ordinary}{m}})][\num{\OLId{latin-ordinary}{n}+\OLId{latin-ordinary}{m}}]$؛
فيلزم $\Th{Q}\Proves\eq[(\num{\OLId{latin-ordinary}{n}}+\num{\OLId{latin-ordinary}{m}}')][\num{\OLId{latin-ordinary}{n}+\OLId{latin-ordinary}{m}}']$.
\end{proof}

\begin{proof}[برهان \olref{prop:rep-add}]
نختار لتمثيل الجمع الـ!!{formula}~$!A_\Add(\OLId{latin-ordinary}{x}_\OLMathNumeral{0},\OLId{latin-ordinary}{x}_\OLMathNumeral{1},\OLId{latin-ordinary}{y})$؛ فتمثيلها
لـ$\Add$ يكون بأخذها $\eq[\OLId{latin-ordinary}{y}][(\OLId{latin-ordinary}{x}_\OLMathNumeral{0}+\OLId{latin-ordinary}{x}_\OLMathNumeral{1})]$.
أما الشرط الأول فهو أنه، متى كان $\Add(\OLId{latin-ordinary}{n},\OLId{latin-ordinary}{m})=\OLId{latin-ordinary}{k}$، لزم
$\Th{Q}\Proves !A_\Add(\num{\OLId{latin-ordinary}{n}},\num{\OLId{latin-ordinary}{m}},\num{\OLId{latin-ordinary}{k}})$، أي
$\Th{Q}\Proves\eq[\num{\OLId{latin-ordinary}{k}}][(\num{\OLId{latin-ordinary}{n}}+\num{\OLId{latin-ordinary}{m}})]$.
وهذا حاصل؛ فإن $\OLId{latin-ordinary}{k}=\OLId{latin-ordinary}{n}+\OLId{latin-ordinary}{m}$، فيكون $\num{\OLId{latin-ordinary}{k}}$ هو $\num{\OLId{latin-ordinary}{n}+\OLId{latin-ordinary}{m}}$ نفسه،
وقد ثبت في \olref{lem:q-proves-add} أن
$\Th{Q}\Proves\eq[(\num{\OLId{latin-ordinary}{n}}+\num{\OLId{latin-ordinary}{m}})][\num{\OLId{latin-ordinary}{n}+\OLId{latin-ordinary}{m}}]$.

وأما الشرط الآخر، فهو أنه متى كان $\Add(\OLId{latin-ordinary}{n},\OLId{latin-ordinary}{m})=\OLId{latin-ordinary}{k}$، فإن
\[
\Th{Q} \Proves \lforall[\OLId{latin-ordinary}{y}][(!A_\Add(\num{\OLId{latin-ordinary}{n}}, \num{\OLId{latin-ordinary}{m}}, \OLId{latin-ordinary}{y}) \lif
  \eq[\OLId{latin-ordinary}{y}][\num{\OLId{latin-ordinary}{k}}])].
\]
لنفرض، باستعمال تناظر المساواة، $\eq[(\num{\OLId{latin-ordinary}{n}}+\num{\OLId{latin-ordinary}{m}})][\OLId{latin-ordinary}{y}]$. وقد ثبت
\[
\Th{Q} \Proves \eq[(\num{\OLId{latin-ordinary}{n}}+\num{\OLId{latin-ordinary}{m}})][\num{\OLId{latin-ordinary}{n}+\OLId{latin-ordinary}{m}}],
\]
فنضع في الطرف الأيسر $\num{\OLId{latin-ordinary}{n}+\OLId{latin-ordinary}{m}}$، فنحصل على
$\eq[\num{\OLId{latin-ordinary}{n}+\OLId{latin-ordinary}{m}}][\OLId{latin-ordinary}{y}]$ لأي~$\OLId{latin-ordinary}{y}$؛ وبالتناظر والتعميم يتم المطلوب.
\end{proof}

\begin{prop}
\ollabel{prop:rep-mult}
ولدالة الضرب
$\Mult(\OLId{latin-ordinary}{x}_\OLMathNumeral{0},\OLId{latin-ordinary}{x}_\OLMathNumeral{1})=\OLId{latin-ordinary}{x}_\OLMathNumeral{0}\cdot \OLId{latin-ordinary}{x}_\OLMathNumeral{1}$ تمثيل في~$\Th{Q}$ بالصيغة:
\[
  !A_{\Mult}(\OLId{latin-ordinary}{x}_\OLMathNumeral{0}, \OLId{latin-ordinary}{x}_\OLMathNumeral{1}, \OLId{latin-ordinary}{y}) \ident \OLId{latin-ordinary}{y} = (\OLId{latin-ordinary}{x}_\OLMathNumeral{0} \times \OLId{latin-ordinary}{x}_\OLMathNumeral{1}).
\]
\end{prop}

\begin{proof} يُترك إثباتها تمرينًا. \end{proof}

\begin{lem}
\ollabel{lem:q-proves-mult}
$\Th{Q} \Proves \eq[(\num{\OLId{latin-ordinary}{n}} \times \num{\OLId{latin-ordinary}{m}})][\num{\OLId{latin-ordinary}{n} \cdot \OLId{latin-ordinary}{m}}]$
\end{lem}

\begin{proof} يُترك إثباتها تمرينًا. \end{proof}

\begin{prob}
ضع برهانًا لـ\olref[inc][req][bre]{lem:q-proves-mult}.
\end{prob}

\begin{prob}
استخرج من \olref[inc][req][bre]{lem:q-proves-mult} برهانًا على
\olref[inc][req][bre]{prop:rep-mult}.
\end{prob}

\begin{explain}
  لا بد من حفظ الفرق بين الرمزين: $\times$ رمز دالة في لغة الحساب،
  وأما $\cdot$ فهو عملية الضرب المألوفة في الأعداد.
  لذلك يوضع $\cdot$ بين تعبيرين عن عددين، نحو $\OLId{latin-ordinary}{m}\cdot \OLId{latin-ordinary}{n}$؛
  وأما $\times$ فموضعه بين حدين من لغة الحساب، نحو
  $(\num{\OLId{latin-ordinary}{m}}\times\num{\OLId{latin-ordinary}{n}})$. ويزيد الأمر التباسًا اشتراك الرمز $+$
  بين !!{function} وعملية الجمع. فإذا وقع بين حدين،
  نحو $(\num{\OLId{latin-ordinary}{n}}+\num{\OLId{latin-ordinary}{m}})$، فهو !!{function} ثنائية الموضع في لغة الحساب؛
  وإذا وقع بين عددين، نحو $\OLId{latin-ordinary}{n}+\OLId{latin-ordinary}{m}$، فهو عملية الجمع.
  وهذا هو المعنى أيضًا في $\num{\OLId{latin-ordinary}{n}+\OLId{latin-ordinary}{m}}$: فالعبارة اسم العدد القياسي الذي
  يقابل العدد~$\OLId{latin-ordinary}{n}+\OLId{latin-ordinary}{m}$.
\end{explain}

\end{document}
